\section{The Transformer Encoder}\label{sec:bg:transformer}

The transformer architecture replaces recurrence with a stack of layers built around \emph{self-attention}, so that every position in a sequence can directly attend to every other position in parallel \parencite{vaswani2017attention}. One understanding of this is a weighted fully connected graph.
Encoder-only stacks, as in \gls{bert} \parencite{devlin2019bert}, map an input sequence to contextualised token representations of the same length and are a natural choice when the downstream task is representation learning or clustering over tokens rather than autoregressive generation.

\paragraph{Scaled dot-product attention.}
Attention maps three matrices of \emph{queries} $\mat{Q}$, \emph{keys} $\mat{K}$, and \emph{values} $\mat{V}$ to an output matrix of the same leading dimension as $\mat{Q}$.
Each row of $\mat{Q}$ scores all rows of $\mat{K}$ through inner products; the scores are scaled, normalised with a row-wise $\softmax$, and used to form a convex combination of the rows of $\mat{V}$.
With $d_k$ the dimension of each query and key vector (column width of $\mat{Q}$ and $\mat{K}$ after the usual transpose convention), the mapping is
\begin{equation}\label{eq:bg:attention}
  \mathrm{Attention}(\mat{Q}, \mat{K}, \mat{V})
  =
  \softmax\!\left(\frac{\mat{Q}\mat{K}^{\mathsf{T}}}{\sqrt{d_k}}\right)\mat{V}.
\end{equation}
The scaling by $\sqrt{d_k}$ keeps the inner products from growing too large in magnitude as $d_k$ increases, which would otherwise drive the $\softmax$ into near one-hot rows and yield vanishing gradients \parencite{vaswani2017attention}.
\cref{eq:bg:attention} is reused later as the core attention primitive inside the encoder.

\paragraph{Self-attention and \gls{mha}.}
In \emph{self-attention}, $\mat{Q}$, $\mat{K}$, and $\mat{V}$ are all linear projections of the same input sequence layout, so a token's representation is updated by aggregating information from the full sequence.
\Gls{mha} runs $h$ attention mechanisms in parallel on different learned subspaces, concatenates the head outputs, and applies one more linear map to restore the model width \parencite{vaswani2017attention}.
The heads specialise to different relational patterns while sharing the same residual and normalisation wrapper as a single conceptual operator.

\paragraph{Positional information.}
Self-attention is permutation-equivariant in its inputs: reordering the rows of $\mat{Q}$, $\mat{K}$, and $\mat{V}$ in the same way reorders the output rows accordingly, but does not otherwise encode order.
Transformer encoders therefore inject a positional signal. The original architecture adds fixed sinusoidal features or learned embeddings to each token vector before the first layer \parencite{vaswani2017attention,devlin2019bert}.
Those encodings are absolute. They tag a token with its index, not with its distance to other tokens.

\paragraph{Rotary positional encoding.}
\textcite{su2024roformer} introduce \gls{rope}. The query and key of each token are rotated by a phase that depends on a scalar coordinate $p$, leaving the token embeddings unchanged.
In the original formulation $p_i$ is the token index $i$. Write $\vect{q}_i = \mat{W}_q\vect{h}_i$ and $\vect{k}_i = \mat{W}_k\vect{h}_i$ for the usual per-head projections of a token representation $\vect{h}_i$. The rotated vectors are functions $f_q(\vect{h}_i, p_i)$ and $f_k(\vect{h}_j, p_j)$ of the hidden representation and the coordinate. Their inner product is required to depend on the two coordinates only through their difference,
\begin{equation}\label{eq:bg:rope_goal}
  \bigl\langle f_q(\vect{h}_i, p_i),\, f_k(\vect{h}_j, p_j) \bigr\rangle
  = g(\vect{h}_i, \vect{h}_j, p_j - p_i).
\end{equation}

Assume $d_k$ is even. The $d_k$ coordinates of a head are grouped into $d_k/2$ two-dimensional planes. Plane $m$ is rotated by the angle $p\omega_m$, where $\omega_m = b^{-2(m-1)/d_k}$ and $b=10^4$ is a fixed base.
Because $p$ increases by one between consecutive tokens, the first plane rotates by one radian per token.
The base $b$ sets the other end of the ladder: the last plane rotates by about $p/b$ radians and completes a full turn only after $p$ reaches $2\pi b$.
Fast-rotating planes resolve small gaps in $p$. Slow-rotating planes resolve large ones. The same geometric progression of wavelengths appears in the sinusoidal encoding of \textcite{vaswani2017attention}.
Each plane uses the rotation
\begin{equation}\label{eq:bg:rope_2d}
  \mat{\Omega}(\phi) =
  \begin{bmatrix}
    \cos\phi & -\sin\phi \\
    \sin\phi & \phantom{-}\cos\phi
  \end{bmatrix},
\end{equation}
and the head-wide map is the block-diagonal matrix
\begin{equation}\label{eq:bg:rope_block}
  \mat{R}(p)
  = \diag\!\bigl(
      \mat{\Omega}(p\omega_1),\,
      \mat{\Omega}(p\omega_2),\,
      \ldots,\,
      \mat{\Omega}(p\omega_{d_k/2})
    \bigr).
\end{equation}
Applying $\mat{R}(p)$ after the projections gives
\begin{equation}\label{eq:bg:rope_map}
  \tilde{\vect{q}}_i = \mat{R}(p_i)\,\vect{q}_i,
  \qquad
  \tilde{\vect{k}}_j = \mat{R}(p_j)\,\vect{k}_j.
\end{equation}

Planar rotations compose by subtracting angles, $\mat{\Omega}(\phi)^{\top}\mat{\Omega}(\psi) = \mat{\Omega}(\psi - \phi)$. Block-wise this is $\mat{R}(p_i)^{\top}\mat{R}(p_j) = \mat{R}(p_j - p_i)$, so
\begin{equation}\label{eq:bg:rope}
  \bigl\langle \tilde{\vect{q}}_i,\, \tilde{\vect{k}}_j \bigr\rangle
    = \vect{q}_i^{\top}\,\mat{R}(p_j - p_i)\,\vect{k}_j.
\end{equation}
Each token is rotated by its own coordinate, yet the logit sees only the difference. Setting $\mat{R}(p) = \mat{I}$ recovers the content-only inner product. \Cref{fig:bg:rope} shows the geometry on a single plane.

The construction applies to any real scalar $p$, not only to integer indices. When the token features already carry timing, a separate positional encoding is sometimes omitted altogether. The method chapter returns to this for \gls{pdw} windows: a baseline that omits a separate encoding, and a rotary encoding that uses physical coordinates in place of token indices.

\begin{figure}[tbp]
  \centering
  % =====================================================================
% Geometry of RoPE on one 2D plane of an attention head.
% Left: absolute phase applied to the projected query.
% Right: content angle phi (dashed) vs. RoPE-modified enclosed angle.
% Labels sit outside the unit circle; angles are large enough to read.
% Loaded from sections/02_background/07_transformer_encoder.tex via \input.
% =====================================================================
\begin{tikzpicture}[
    font=\small,
    >=Latex,
    axis/.style={->, gray!55, thin},
    guide/.style={gray!30, thin},
    qraw/.style={->, thick, densely dashed, color=black!50},
    qvec/.style={->, very thick, color=blue!70!black},
    kraw/.style={->, thick, densely dashed, color=red!45!black},
    kvec/.style={->, very thick, color=red!70!black},
    angarc/.style={->, thin, color=black!70},
  ]

  % =================== LEFT PANEL =======================================
  \begin{scope}[xshift=0mm]
    \def\Rad{18mm}
    \node[font=\footnotesize\bfseries, anchor=south] at (3mm, 30mm)
      {Apply $\mat{R}(p)$ to the projected query};

    \draw[axis] (-20mm, 0) -- (24mm, 0);
    \draw[axis] (0, -11mm) -- (0, 24mm);
    \draw[guide] (0, 0) circle (\Rad);

    % ~35° absolute phase
    \def\qbase{15}
    \def\arot{50}

    \draw[qraw] (0, 0) -- (\qbase:\Rad);
    \node[color=black!55, font=\scriptsize, anchor=west, inner sep=1.5pt]
      at (\qbase:1.08*\Rad) {$\vect{q}=\mat{W}_q\vect{h}$};

    \draw[qvec] (0, 0) -- (\arot:\Rad);
    \node[color=blue!70!black, font=\scriptsize, anchor=south west, inner sep=1.5pt]
      at (\arot:1.08*\Rad) {$\tilde{\vect{q}}=\mat{R}(p)\,\vect{q}$};

    \draw[angarc] (\qbase:7.5mm) arc (\qbase:\arot:7.5mm);
    \draw[gray!55, thin]
      ({(\qbase+\arot)/2}:7.5mm) -- (1.35*\Rad, 0.55*\Rad);
    \node[font=\scriptsize, anchor=west] at (1.38*\Rad, 0.55*\Rad)
      {$p\omega_m$};

    \node[font=\scriptsize, anchor=north, align=center, text=black!60]
      at (2mm, -12mm)
      {projection first, then a planar rotation\\[-1pt]
       by the coordinate $p$};
  \end{scope}

  % =================== RIGHT PANEL ======================================
  \begin{scope}[xshift=76mm]
    \def\Rad{18mm}
    \node[font=\footnotesize\bfseries, anchor=south] at (3mm, 30mm)
      {Inner product sees only $p_j-p_i$};

    \draw[axis] (-20mm, 0) -- (24mm, 0);
    \draw[axis] (0, -11mm) -- (0, 24mm);
    \draw[guide] (0, 0) circle (\Rad);

    % content φ = 20°; relative phase = 28°; vectors well separated
    \def\qrawang{10}
    \def\krawang{30}
    \def\qang{45}
    \def\kang{93}

    \draw[qraw] (0, 0) -- (\qrawang:\Rad);
    \draw[kraw] (0, 0) -- (\krawang:\Rad);
    \node[color=black!50, font=\scriptsize, anchor=west, inner sep=1.5pt]
      at (\qrawang:1.08*\Rad) {$\vect{q}_i$};
    \node[color=red!45!black, font=\scriptsize, anchor=south west, inner sep=1.5pt]
      at (\krawang:1.08*\Rad) {$\vect{k}_j$};

    \draw[angarc, color=black!45] (\qrawang:5mm) arc (\qrawang:\krawang:5mm);
    \node[font=\scriptsize, color=black!55, anchor=west]
      at ({(\qrawang+\krawang)/2}:6.2mm) {$\varphi$};

    \draw[qvec] (0, 0) -- (\qang:\Rad);
    \draw[kvec] (0, 0) -- (\kang:\Rad);
    \node[color=blue!70!black, font=\scriptsize, anchor=south west, inner sep=1.5pt]
      at (\qang:1.08*\Rad) {$\tilde{\vect{q}}_i$};
    \node[color=red!70!black, font=\scriptsize, anchor=south, inner sep=1.5pt]
      at (\kang:1.08*\Rad) {$\tilde{\vect{k}}_j$};

    % Outer arc; label radially outward from the arc midpoint (avoids crossing q̃_i)
    \draw[angarc] (\qang:11mm) arc (\qang:\kang:11mm);
    \pgfmathsetmacro{\relang}{(\qang+\kang)/2}
    \draw[gray!55, thin] (\relang:11mm) -- (\relang:1.38*\Rad);
    \node[font=\scriptsize, anchor=south west, inner sep=1.5pt]
      at (\relang:1.4*\Rad) {$\varphi+(p_j-p_i)\omega_m$};

    \node[font=\scriptsize, anchor=north, align=center, text=black!60]
      at (2mm, -12mm)
      {$\langle\tilde{\vect{q}}_i,\tilde{\vect{k}}_j\rangle
        =\vect{q}_i^{\top}\mat{R}(p_j-p_i)\,\vect{k}_j$};
  \end{scope}

\end{tikzpicture}

  \caption{Geometry of \gls{rope} on one two-dimensional plane of a head. Left: the projected query $\vect{q}=\mat{W}_q\vect{h}$ (dashed) is rotated by $\mat{R}(p)$ from \cref{eq:bg:rope_block}, which adds the phase $p\omega_m$. Right: the angle between the rotated query and key equals the content angle $\varphi$ plus the relative phase $(p_j-p_i)\omega_m$. The inner product depends on the coordinate difference alone (\cref{eq:bg:rope}).}
  \label{fig:bg:rope}
\end{figure}

\paragraph{Encoder layer.}
A standard transformer \emph{encoder layer} alternates two sub-layers: \gls{mha} followed by a position-wise \gls{ffn} applied independently at every position.
Each sub-layer uses a residual connection and layer normalisation \parencite{vaswani2017attention}.
The \gls{ffn} is typically a two-layer \gls{mlp} with a wide inner dimension and a nonlinear activation, acting as the per-token nonlinearity while \gls{mha} handles cross-token mixing.
Stacking $N$ such layers yields a deep encoder that maps an input sequence $(\vect{x}_1,\ldots,\vect{x}_L)$ to contextualised representations $(\vect{h}^{(N)}_1,\ldots,\vect{h}^{(N)}_L)$, which downstream heads can pool or read out for sequence-level or token-level predictions.

\paragraph{Computational cost.}
In the usual dense implementation, each self-attention block materialises an $L \times L$ matrix of token--token scores before the softmax, so time and activation memory for that block scale as $\mathcal{O}(L^2)$ in the sequence length $L$ at fixed embedding width, up to constants from the number of heads and projection ranks \parencite{vaswani2017attention}.
Repeating the block across depth multiplies that cost by the number of layers, which motivates hard caps on $L$, sparse or low-rank attention variants, or chunked processing when inputs are long; the method chapter states the windowing choice adopted here.
